\section{Finite reconstruction from one unknown launch law}
\label{sec:single-law-finite}
The angular identity also has a finite statistical realization.  Its
regularization must preserve positivity: a pointwise numerical second
derivative of a measured field would neither supply the necessary
support test nor admit a uniform bias estimate for arbitrary $L^1$
densities.  We instead differentiate a known angular test function and
average the spatial command centers with positive weights.  Both
averages are implemented by finite rational menus of ordinary forward
commands.

In this section $s=6+\beta$, and the obstacle class has the bounds of
\eqref{eq:priors}, including uniformly positive upper and lower
curvature radii.  All widths are at least a known $w_*>0$.  The one
unknown footprint and density satisfy
\begin{equation}\label{eq:single-finite-priors}
 \begin{gathered}
 \diam A\le\Delta<\min(d_0,w_*),\qquad
 \|p_A^\circ\|_{C^{6,\beta}}\le M_A,\\
 0<r_A^-\le p_A^\circ+(p_A^\circ)''\le r_A^+,
 \qquad j(z)\ge b_0\dist(z,\partial A)^\gamma
       \quad\hbox{for a.e. }z\in A,
 \end{gathered}
\end{equation}
where $b_0>0$ and $\gamma\ge0$ are fixed.  A uniform inball bound
could equivalently be added to these priors; the displayed lower
curvature radius already supplies one.  No upper bound, smoothness,
translation modulus, or evaluations of $j$ are available.  Put
\[
 a_0=s(A),\qquad A_0=A-a_0,\qquad
 \mathcal O_0=\mathcal O-a_0,
 \qquad \kappa_*:=\min(d_0,w_*)-\Delta>0.
\]
Here $s(A)$ denotes the Steiner point.  It is an unknown used to state
the identifiable output, not a supplied coordinate.  The experiment
uses only the means \eqref{eq:single-raw-response} and the individual
forward bits whose means they are.  All nominal centers and displacement
vectors selected below are rational.

\subsection{A positive weak response and its finite estimator}
Let $\chi(u)=c_\chi(1-u^2)^4$ for $|u|\le1$ and zero otherwise,
where $c_\chi$ normalizes its integral to one.  This constant is
rational.  Define
\begin{equation}\label{eq:single-positive-kernels}
 \kappa_r(u)=r^{-2}\chi(u_1/r)\chi(u_2/r),\qquad
 K_b(v)=b^{-1}\chi(v/b),\qquad \lambda_b=K_b''+K_b.
\end{equation}
The kernels have compact support and enough endpoint vanishing for
twice integrating by parts.  They may also be approximated by smooth
test functions in that identity.  Their bounds include
\begin{equation}\label{eq:single-kernel-bounds}
 \|\kappa_r\|_\infty\le Cr^{-2},\quad
 \|\nabla\kappa_r\|_\infty\le Cr^{-3},\quad
 \|\lambda_b\|_1\le Cb^{-2},\quad
 \|\lambda_b'\|_1\le Cb^{-3},\quad
 \int\lambda_b=1.
\end{equation}
Angles are integrated on a real interval around the chosen direction;
the physical responses are extended periodically.  Thus no choice of
a branch for an angular coordinate enters the construction.

Let $S_\theta$ be the positive angular measure in
\eqref{eq:single-angular-field}.  At a spatial target $x$ and angular
target $\alpha$ define the nonnegative function
\begin{align}
 U_{r,b}(x,\alpha)
  &=\int_{\R}\!K_b(v)\int_{\R^2}\!\kappa_r(u)
                      S_{\alpha+v}(x-u)\dd u\dd v
                       \label{eq:single-positive-response}\\
  &=\int_{\R}\!\lambda_b(v)\int_{\R^2}\!\kappa_r(u)
                      q_{\alpha+v}(x-u)\dd u\dd v.
                       \label{eq:single-weak-response}
\end{align}
The first formula defines a continuous nonnegative spatial function;
the second follows from Lemma~\ref{lem:single-angular-resolution}.
These are convolutions of locally finite measures.  In particular,
neither expression requires pointwise values of $j$.

\begin{lemma}[Finite sampling of the weak angular response]
\label{lem:single-weak-sampling}
Fix a bounded nominal aperture and all the stated priors.  For
$0<r,b,t<1$ sufficiently small, let $\ell$ be a spatial quadrature
mesh, $d_\theta$ an angular quadrature mesh, and $\zeta$ a bound on
the rounding of each displacement vector, with
\[
 \ell\le c\min(r,t),\qquad d_\theta\le cb,\qquad \zeta\le ct.
\]
There is a finite estimator $\widehat U(x,\alpha)$ using only forward
collision bits, whose deterministic bias relative to
$U_{r,b}(x,\alpha)$ is at most
\begin{equation}\label{eq:single-weak-bias}
 C\left\{tr^{-3}b^{-2}
       +\frac{\ell r^{-2}b^{-2}}t
       +d_\theta r^{-2}b^{-3}
       +\frac{\zeta r^{-2}b^{-2}}t\right\}.
\end{equation}
For any $a>0$ and $0<\varepsilon<1$, an additional statistical error
at most $a$ has probability at least $1-\varepsilon$ after
\begin{equation}\label{eq:single-weak-cost}
 n\le C\left(t^{-1}r^{-2}b^{-4}a^{-2}
                  +t^{-1}b^{-2}a^{-1}\right)
                    \log\frac{C}{\varepsilon}
\end{equation}
attempted bits, with an integer ceiling understood.  The constants
are uniform over the unknown density.  A fresh rational center and
vector selection is charged as a command occurrence, even when it
is chosen by randomizing a nominal integral.
\end{lemma}
\begin{proof}
First replace $q_\theta$ in \eqref{eq:single-weak-response} by
$t^{-1}F_{t,\theta}$.  Apply
\eqref{eq:single-weak-flight-error} to
$\psi(y)=\kappa_r(x-y)$ and integrate against $|\lambda_b|$.
Separation and the fixed diameter and aperture bounds give a uniform
bound on the number and widths of all components meeting these tests.
The first term in \eqref{eq:single-weak-bias} follows.  Integrating
over $j$ costs only its total mass one.  Thus an unobserved modulus
of continuity of $j$ has not been used in this replacement.

We give the finite quadrature details because rounding pointwise
values of an arbitrary density would not give this conclusion.  Put
\[
 w_u=\ell^2\kappa_r(u),\quad u\in\ell\Z^2,\qquad
 Z=\sum_u w_u,
 \qquad v_k=d_\theta\lambda_b(kd_\theta),\quad B=\sum_k|v_k|.
\]
Only the supports of the kernels are included.  For small fixed mesh
ratios, $1/2\le Z\le2$ and $B\le Cb^{-2}$.  For dyadic $r,b,\ell$
and $d_\theta$, all these weights are rational and are specified by
the fixed polynomials in \eqref{eq:single-positive-kernels}.

For a fixed launch displacement $z$ and fixed command $a$, the
function of the spatial variable being quadratured is a smooth weight
times the indicator of a collision strip.  For $|a|<d_0$, such a
strip is the difference between a convex swept body and its original
convex body.  In the relevant bounded region their total perimeter
is uniformly bounded.  The union of square mesh cells of side $\ell$
meeting these boundaries has area at most $C\ell$; this follows by
covering each convex boundary by its inner and outer parallel bands.
The weight variation on the remaining cells has integral at most
$C\ell r^{-1}$.  Since $\|\kappa_r\|_\infty\le Cr^{-2}$, the total
quadrature error is at most $C\ell r^{-2}$.  The estimate is uniform
in the translation $z$.  Integration against $j(z)\dd z$, division
by $t$ and summation against $|v_k|$ therefore give the second term
in \eqref{eq:single-weak-bias}.

For the angular quadrature write
\[
 A_t(\theta)=\int\kappa_r(u)F_{t,\theta}(x-u)\dd u.
\]
The swept strip has area at most $Ct$, and changing its displacement
vector from $t n_\theta$ to $t n_{\theta'}$ changes its area in
symmetric difference by at most $Ct|\theta-\theta'|$.  To see the
latter estimate, their convex swept bodies have Hausdorff distance
at most $t|n_\theta-n_{\theta'}|$ and uniformly bounded perimeter;
the unchanged solid is subtracted from both.  Consequently
\[
 0\le A_t(\theta)\le Ct r^{-2},\qquad
 |A_t(\theta)-A_t(\theta')|
                    \le Ct r^{-2}|\theta-\theta'|.
\]
These estimates remain uniform after integration over $j$.  The
one-dimensional quadrature error for $\lambda_b A_t$ is bounded by
its total variation times $d_\theta$.  Equation
\eqref{eq:single-kernel-bounds}, followed by division by $t$, gives
$Cd_\theta r^{-2}b^{-3}$.  Apply this argument to the continuous
spatial average; the spatial quadrature error at the finitely many
angular nodes has already been bounded separately.

Round $t n_{\alpha+kd_\theta}$ to a rational vector $a_k$ at distance
at most $\zeta$.  The same swept-body area estimate gives an error
$C\zeta r^{-2}$ in each continuous spatial mean.  The spatial
quadrature estimate is valid for the rounded vectors as well, so
this gives the final term of \eqref{eq:single-weak-bias}.

For each attempt select $u$ with probability $w_u/Z$, and select $k$
independently with probability $|v_k|/B$.  Attempt the ordinary
forward command at nominal center $x-u$ with displacement $a_k$, and
write its collision bit as $Y$.  The signed random variable
\begin{equation}\label{eq:single-weighted-observation}
                   W=\frac{ZB}{t}\operatorname{sgn}(v_k)Y
\end{equation}
has expectation exactly equal to the finite quadrature just described.
Zero weights can be omitted.  This selection does not modify the
stationary launch law: the random choices select known nominal
commands, each followed by a fresh draw from the same unknown $j$.

Uniformly at every angular node, the mean collision probability under
the spatial proposal is at most
$C(t+\ell+\zeta)r^{-2}\le Ct r^{-2}$.  It follows that
\begin{equation}\label{eq:single-weak-variance}
 |W|\le Ct^{-1}b^{-2},\qquad
 \operatorname{Var}(W)\le\E W^2\le Ct^{-1}r^{-2}b^{-4}.
\end{equation}
Bernstein's inequality for the average of independent copies of $W$
proves \eqref{eq:single-weak-cost}.  All observations, including zeros
from solid starts, enter this average.  The selection probabilities
and multipliers are known rational numbers; no response mean or
unknown density is an input to them.
\end{proof}

\begin{lemma}[Separated footprint copies from a finite threshold field]
\label{lem:single-finite-copies}
Let $0<e<e_0$ and choose $r,b$ comparable to sufficiently small fixed
multiples of $e$.  At each angular target $\alpha$ the function
$U_{r,b}(\cdot,\alpha)$ is supported in the union of the $Ce$-parallel
bodies of
\begin{equation}\label{eq:single-target-copies}
                 K_{C,\sigma}=c_C(\alpha+\sigma\pi/2)-A,
                         \qquad \sigma\in\{-1,1\}.
\end{equation}
These parallel bodies remain separated by at least $\kappa_*/2$.
There is a known $c_1>0$ such that
\begin{equation}\label{eq:single-interior-signal}
 U_{r,b}(x,\alpha)\ge c_1e^\gamma
            \quad\hbox{if }x\in(K_{C,\sigma})_{-e}.
\end{equation}
Suppose its values on a spatial mesh of width $c e$ in a protected
aperture are estimated with error at most $c_1e^\gamma/8$.
Thresholding, a fixed-distance graph, and convex hulls then recover
every complete protected copy in Hausdorff distance at most $Ce$,
with exactly one recovered cluster per copy.  Their Steiner points
estimate $c_C(\alpha+\sigma\pi/2)-a_0$ with error at most $Ce$.
\end{lemma}
\begin{proof}
Write $\phi=\alpha+\sigma\pi/2$.  The radius bound gives
$|c_C(\phi+v)-c_C(\phi)|\le r_C^+|v|$.  Each term in the first
formula \eqref{eq:single-positive-response} is therefore supported
in $K_{C,\sigma}+(r_C^+b+\sqrt2r)\overline B$.  Choose these two
widths to have sum at most $e/4$.  The exact separation of the copies
is at least $\kappa_*$, proving the support and gap assertions.

If $x\in(K_{C,\sigma})_{-e}$, then every displacement
$c_C(\phi+v)-x+u$ occurring in this term lies inside $A$ at depth
at least $3e/4$.  The lower bound in
\eqref{eq:single-finite-priors}, the lower obstacle curvature radius,
and the unit integrals of both positive kernels give
\eqref{eq:single-interior-signal}.  The other terms are nonnegative.
This argument integrates the almost-everywhere density inequality;
it requires no pointwise choice of its representative.

Retain grid nodes whose estimated value exceeds $c_1e^\gamma/4$.
Every node at depth $e$ is retained, and no node outside the
$e/4$-parallel bodies is retained.  Give two retained nodes an edge
when their distance is less than $\kappa_*/4$.  Distinct copies cannot
be joined for small $e$.  Within a copy, the nodes in a slightly deeper
erosion form a connected set for this fixed edge length: segments
joining points of the convex erosion can be followed by a grid chain
with error $O(e)$.  Every retained node lies within $Ce$ of that
erosion and hence of a retained deep node.  To justify this last
assertion uniformly up to the boundary, the footprint's lower
curvature radius implies that its inner parallel body at distance
$u<r_A^-$ has support $p_A-u$, so its Hausdorff distance from $A$
is $u$.  Alternatively its fixed inball gives the same bound by a
homothety towards an inball center.  Thus each complete copy yields
exactly one graph component.

The hull of its retained nodes lies in $K_{C,\sigma}+Ce\overline B$
and contains a grid approximation of $(K_{C,\sigma})_{-2e}$.
Its support consequently differs from that of $K_{C,\sigma}$ by
at most $Ce$.  Steiner centering is Lipschitz for support supremum
error, with constant two in the plane, proving the final assertion.
\end{proof}

\begin{theorem}[Finite forward reconstruction at one unknown setting]
\label{thm:single-law-finite}
Fix the priors \eqref{eq:single-finite-priors} and the bounded periodic
class $\mathcal B_\eta$.  At a single unknown stationary launch law,
a finite controller using only forward collision bits reconstructs
$A_0$ and a smooth strictly convex periodic table approximating
$\mathcal O_0$.  For all sufficiently small $\nu>0$ and
$0<\delta<1/4$, its footprint $C^2$ support error and the geometric
loss of the table are at most $C\nu$, and its primitive discrete data
are correct, with probability at least $1-\delta$.

Put
\begin{equation}\label{eq:single-finite-exponent}
                 Q_\gamma=\frac{(3\gamma+27/2)s}{s-2}.
\end{equation}
The controller has the deterministic attempted-bit bound
\begin{equation}\label{eq:single-finite-resources}
 N_\nu\le C\nu^{-Q_\gamma}\log\frac{C}{\nu\delta},
             \qquad S_\nu\le J_\nu\le N_\nu,
\end{equation}
where $J_\nu$ counts center-and-vector-labelled command occurrences
and $S_\nu$ counts distinct nominal sites.  The nominal centers remain in a
fixed prior-dependent aperture.  If a known coarse bound on $|a_0|$ is
also imposed, all actual launches and attempted flights lie in a
prior-bounded physical aperture; the bound does not supply the origin.
The actual starts lie in $W+A+O(r)\overline B$, and their flights
in its $t$-parallel body.  Its directional commands have lengths
of order
\begin{equation}\label{eq:single-finite-precision}
 \nu^{(\gamma+5)s/(s-2)},
 \qquad\hbox{and coordinate mesh of order }
             \nu^{(2\gamma+9)s/(s-2)}.
\end{equation}
The total command description is at most
$C J_\nu\log(C/(\nu\delta))$ bits, for fixed computable priors and
certified dyadic choices of these scales.  Numerical construction of
this finite menu and the geometric output is separate from the
observation count.  No sharpness of the exponent $Q_\gamma$ is asserted.
The fixed-known-disk minimax theorem concerns its own experiment and
retains its stated exponent.

The footprint translation $a_0$ is not estimated separately, as it
belongs to the common translation ambiguity of the experiment.  There
is one setting, no homothety control, no reciprocal preparation, and
no supplied density or footprint evaluations.  The theorem estimates
the geometry; it makes no strong-norm finite density estimate under
an unrestricted $L^1$ density class.
\end{theorem}
\begin{proof}
Choose a dyadic $e$ comparable to a sufficiently small multiple of
$\nu^{s/(s-2)}$.  In Lemma~\ref{lem:single-weak-sampling} take dyadic
scales comparable to
\begin{equation}\label{eq:single-finite-scales}
 r\asymp e,\quad b\asymp e,\quad
 t\asymp e^{\gamma+5},\quad
 \ell\asymp\zeta\asymp e^{2\gamma+9},\quad
 d_\theta\asymp e^{\gamma+5}.
\end{equation}
The constants in $r,b$ first meet the support reserves of
Lemma~\ref{lem:single-finite-copies}; the constant in $t$ and then
those in $\ell,d_\theta,\zeta$ make the four bias terms sum to at
most $c_1e^\gamma/16$.  These choices satisfy all mesh and length
restrictions.  With additional statistical error
$a=c_1e^\gamma/16$, \eqref{eq:single-weak-cost} becomes
\begin{equation}\label{eq:single-per-target-count}
                  n\le C e^{-(3\gamma+11)}
                                   \log(C/\varepsilon).
\end{equation}
Indeed the two Bernstein powers are $3\gamma+11$ and $2\gamma+7$,
and the first dominates.  The spatial boundary quadrature, rather
than an assumed continuity modulus for $j$, justifies the much finer
rational command mesh in \eqref{eq:single-finite-scales}.

We specify the aperture and discard rule before using these estimates.
Choose a fixed radius $R$ that contains the entire protected patch
needed by Lemma~\ref{lem:locking}, in the frame of $\mathcal O_0$.
A presentation lattice with the fixed bounds supplies representatives
in a prior-bounded box after any common translation, so no coordinate
of $a_0$ is needed for this choice.  Take a square $W$ whose interior
contains the ball of radius
$R+3D_0+5\Delta+10$ and its fixed response collars.  Sampling takes
place at the spatial target grid in $W$; the nominal centers used by
the quadrature extend this only by $\sqrt2r$.

At each angular target form the graph of Lemma~\ref{lem:single-finite-copies}.
Keep only graph components with a node at distance greater than
$2\Delta+1$ from $\partial W$.  Any such node is within $Ce$ of its
true copy, whose diameter is at most $\Delta$.  Its entire copy,
including the required $Ce$ margin, is therefore contained in $W$.
Consequently its hull and Steiner point satisfy that lemma; a
truncated exterior copy cannot be kept.  Conversely all copies of
obstacles meeting the protected region occur with the required
margin, at every angle.  This gives a finite complete-component
cutoff without an occupation query or a geometric seed supplied to
the controller.

Use angular targets in a rational mesh of maximum gap
$d=c\sqrt e$ around the circle.  There are at most $Ce^{-1/2}$ of
them, and at most $Ce^{-2}$ spatial targets at each.  Rational
approximations to $2\pi$ with a reserved margin give such a mesh;
the actual vectors at all quadrature nodes are rounded with the
certified error $\zeta$.  Apply a union bound over these $Ce^{-5/2}$
targets, setting their failure allowances to
$c\delta e^{5/2}$.  Equations \eqref{eq:single-per-target-count}
and \eqref{eq:single-finite-scales} give the attempted-bit bound in
\eqref{eq:single-finite-resources}.

On the resulting common event, let $\mathcal Z$ be the union of the
estimated Steiner points of retained copies.  Each point lies within
$Ce$ of $c_C(\phi)-a_0$ for some obstacle and the corresponding
sampled normal angle.  Join these points when their distance is less
than $d_0/3$.  Points coming from different obstacles cannot be
joined.  For a complete protected obstacle, consecutive angular
samples of its Gauss parametrization are at distance at most
$r_C^+d+Ce<d_0/3$, so all its samples belong to one component.
A cloud component with a point in the protected radius $R+D_0+1$
comes from an obstacle lying wholly in the larger aperture reserve.
All angular copies of that obstacle were retained.  Components from
outside this cutoff are not used as complete obstacles.  Every
obstacle needed for the period tests is present.

For each retained complete cloud component take its convex hull.
This hull estimates the corresponding $C-a_0$ in Hausdorff distance
$Ce$.  Here the angular mesh $\sqrt e$ suffices, despite the absence
of labels distinguishing the two tangent points.  For every normal
angle $\phi$ choose one sampled Gauss angle $\psi$ at distance at
most $d$.  Since $c_C'(u)=r_C(u)\tau_u$,
\begin{equation}\label{eq:single-gauss-support-deficit}
 0\le h_C(\phi)-c_C(\psi)\cdot n_\phi
    \le r_C^+(1-\cos|\phi-\psi|)
    \le\tfrac12r_C^+d^2.
\end{equation}
Thus the exact sampled boundary hull has support error $O(d^2)$;
the error in each measured center adds $Ce$.  The union of samples,
with no across-angle component labels supplied, gives this hull.

Take any complete copy hull from the angular stage.  Reflect its
Steiner-centered support to estimate $A_0$, also in support supremum
error $Ce$.  All true copy hulls are translates of $-A$ and therefore
give the same centered footprint.  Choosing any one of the finitely
many recovered complete copies creates no matching or asymmetry
assumption.

For completeness, these polygonal support estimates yield the claimed
$C^2$ output without treating an arbitrary polygon as a smooth body.
Convolve each support estimate with the signed approximation kernel
of Section~\ref{sec:stationary}, whose nonconstant moments through
degree six vanish, at scale $h=e^{1/s}$.  For a true $C^{6,\beta}$
support $p$ and its estimate $\widetilde p$ of supremum error $Ce$,
Taylor's formula and the $L^1$ derivative bounds of this fixed kernel
give
\[
 \|L_h*\widetilde p-p\|_\infty\le C(e+h^s),\qquad
 \|L_h*\widetilde p-p\|_{C^2}
                         \le C(eh^{-2}+h^{s-2})\le C\nu.
\]
The smoothed supports also have a uniform $C^3$ bound, since their
third derivative error is at most $C(eh^{-3}+h^{s-3})$.  Evaluate them
on an angular mesh of width $b_1$ comparable to a small multiple of
$\min(e^{1/3},\nu)$, to value error $O(b_1^3)$.  The polygonal
input and fixed kernel have explicit Lipschitz bounds for certifying
these finite integrals.  Quadratic local interpolation with the fixed
smooth partition used in Section~\ref{sec:stationary} then adds
supremum error $O(b_1^3)=O(e)$ and $C^2$ error $O(b_1)=O(\nu)$.
It gives a finite smooth support representation.  Steiner centering
changes the same bounds only by a fixed factor.  The known positive lower curvature
radii ensure $\widehat p+\widehat p''>0$ for small $\nu$, so all
outputs are supports of actual smooth strictly convex bodies.

Apply the retained period-locking and periodic assembly argument to
these original component estimates.  The margin transfers to the
frame $\mathcal O-a_0$ by
Lemma~\ref{lem:registration-period-margin}; it is invariant under a
common translation after reduction by presentation periods.  Matched
period vectors and free area have error $O(e)$, the discrete data
are exact on the common event, and the assembled periodic table has
the asserted geometric loss.  This final algebra uses the acquired
forward-data geometry and no reciprocal measurement.

Finally, each sampled center and displacement vector is specified on
the dyadic mesh $\ell\asymp\zeta$, within a fixed aperture and a
fixed displacement bound.  It therefore has
$O(\log(C/\nu))$ coordinate bits.  The kernels, menus and selection
probabilities have fixed polynomial formulas; certified dyadic scale
enclosures and finite rational arithmetic specify them.  Charging
every selected center-and-vector pair separately gives
$S_\nu\le J_\nu\le N_\nu$, with the description bound stated in
the theorem.  This accounting does not confuse an integral's target
with the many physical nominal sites used to estimate it.
\end{proof}

\begin{corollary}[A finite nonperiodic configuration]
\label{cor:single-finite-cloud}
Assume instead that the configuration has between one and a fixed
number $n_0$ of components, with the same geometric and density
priors, and that a known protected aperture contains all expanded
components $C-A$ together with the fixed collars used above.  The
same forward-only controller recovers $A_0$ and all components of
$\mathcal O_0$ in $C^2$ error at most $C\nu$, with probability at
least $1-\delta$ and the resource bounds
\eqref{eq:single-finite-resources}--\eqref{eq:single-finite-precision}.
Neither periodicity nor a nonperiod-patch margin is needed.
\end{corollary}
\begin{proof}
The protected aperture makes every relevant angular support copy
complete.  The two graph constructions and support smoothing in the
proof of Theorem~\ref{thm:single-law-finite} recover every component.
No period test or periodic assembly is performed.  All estimates and
sampling counts preceding that step remain unchanged.
\end{proof}
